\subsection{A bilateral test with observed choices}\label{app:observed-bilateral}
A school may report student preferences while assigning seats by priority. A speed-dating experiment gives us a different check: both sides record whether they would meet again. We used the 16 Columbia sessions in the public Bhargava--Fisman replication archive, joining it to the extended questionnaire by wave and the two participant identifiers. All 7,674 directed decisions agreed. The 474 participants made 644 mutual affirmative choices. Both overall-liking scores were recorded and in range on 640 of those pairs. The experiment allowed multiple contacts; there is no observed one-partner placement. For the counterfactual we treat mutual assent as acceptability, order those partners by the participant's reported overall liking (independently resolving ties), and run unit-capacity DA \citep{fisman2006dating,bhargava2014data}.

We trained on the first eight sessions (222 people) and tested on the last eight (252 people, 2,163 meetings, six distinct event dates). No person or event date crosses the split. A regularized categorical model predicts the full-visibility DA partner or the outside option using pre-event profile features, pairwise age and interest similarity, and self-descriptions weighted by the other person's stated priorities. At quota five, the primary rule ranks everyone with a documented two-sided meeting, before either response or rating is known. Only mutually acceptable displayed edges can lead to a placement. The direct empirical outcome is the fraction of actual mutual yeses that would have received exposure under each display, holding all responses fixed. It measures coverage of recorded choices. It does not identify how people would have responded to actually meeting fewer partners.

\begin{table}[ht]\centering\small
\caption{Observed bilateral choices in eight later sessions. Women propose. Percentages average sessions equally. Recovery is computed among proposers assigned by full-visibility DA, placement among all proposers, and mutual coverage among all recorded mutual yeses.}
\begin{tabular}{@{}lrrr@{}}\toprule
Top-five rule & Full partner on list & Placed by DA & Mutual yes shown\\\midrule
Uniform & 34.94 & 42.79 & 34.72\\
Profile-based partner predictor & 36.32 & 37.15 & 41.43\\
Predictor with assignment anchor & 39.36 & 39.56 & 40.92\\
Age and interest similarity & 42.73 & 48.24 & 42.35\\\bottomrule
\end{tabular}\label{tab:bilateral-real}\end{table}

Against uniform, the predictor gained 1.37 percentage points in DA partner recovery (95\% interval $[-1.73,4.38]$) and lost 5.64 in placement ($[-13.58,0.97]$). The two point estimates have the rental pattern, with uncertainty spanning zero for each. The fraction of actual mutual yeses shown rose by 6.71 points ($[2.13,10.95]$), while the share of potential receivers shown to anyone fell from 99.59\% to 70.53\%. Preference for a real reciprocal contact can therefore be measured separately from placement under single capacity. A simple age-and-interest rule improved both DA point estimates. Reversing the proposer side gave the learned rule a 2.95-point recovery increase and a 3.76-point placement decrease, also imprecise. Intervals use 20,000 paired resamples of the six test dates. Within a date its sessions stay together. The 256 refinements and random display draws per session integrate algorithmic randomness and add no independent data.

A second comparison gives all rules the same accepted graph before computing displays, as in the housing study's known feasibility constraint. In this case 104 of 125 female proposers had at most five acceptable options. Recovery was 96.30\% under learning and 94.28\% under uniform. Placement was 63.89\% and 63.72\%, respectively. At quota three both point estimates again improved ($+1.73$ and $+3.36$). Thus the scarcity of genuinely acceptable partners and what the platform already knows can make a five-option shortlist nearly exhaustive.
